\documentclass[10pt]{article}
\usepackage[T1]{fontenc}
\usepackage[utf8]{inputenc}
\usepackage[french]{babel}
\usepackage{amsmath,amssymb,mathrsfs,mathpazo}
\usepackage{geometry}
\geometry{a4paper,margin=2cm}
\usepackage{graphicx,array,tabularx,tabularray,colortbl,multirow,multicol,booktabs,enumitem,pifont,ulem,fancybox,diagbox,cancel}
\usepackage{pstricks,pstricks-add,pst-plot,pst-tree,pst-node,pst-text,pst-func,pst-math,pst-bspline,pst-3dplot,pst-eucl,pst-solides3d}
\let\pscalculate\undefined
\usepackage{pst-bezier}
\usepackage{tikz,pgfplots,tkz-tab}
\usetikzlibrary{arrows,arrows.meta,patterns,calc,positioning,trees,decorations.markings,shapes.geometric,shapes.misc,shapes.arrows,intersections,decorations.pathmorphing,decorations.pathreplacing,backgrounds,shadows,babel}
\pgfplotsset{compat=1.18}
\usepgfplotslibrary{fillbetween}
\usepackage[np]{numprint}
\usepackage{listings,makecell,pas-tableur,twoopt}
\let\barre\undefined
\usepackage[pstricks]{tablvar}
\definecolor{darkgreen}{RGB}{0,100,0}
\providecommand{\red}{\color{red}}
\providecommand{\blue}{\color{blue}}
\providecommand{\green}{\color{green}}
\providecommand{\black}{\color{black}}
\providecommand{\white}{\color{white}}
\providecommand{\gray}{\color{gray}}
\providecommand{\cyan}{\color{cyan}}
\providecommand{\magenta}{\color{magenta}}
\providecommand{\yellow}{\color{yellow}}
\providecommand{\violet}{\color{violet}}
\providecommand{\orange}{\color{orange}}
\providecommand{\euro}{\texteuro}
\providecommand{\EUR}{\texteuro}
\providecommand{\ieme}{\textsuperscript{e}}
\providecommand{\degres}{\ensuremath{^\circ}}
\providecommand{\degree}{\ensuremath{^\circ}}
\providecommand{\vv}[1]{\overrightarrow{#1}}
\providecommand{\Coord}[1]{#1}
\providecommand{\decosix}{}
\providecommand{\starredbullet}{\ensuremath{\bullet}}
\pagestyle{empty}
\setlength{\parindent}{0pt}
\usepackage{adjustbox}\DeclareUnicodeCharacter{2236}{:}\DeclareUnicodeCharacter{2190}{\ensuremath{\leftarrow}}
\providecommand{\R}{\mathbb{R}}
\providecommand{\N}{\mathbb{N}}
\providecommand{\Z}{\mathbb{Z}}
\providecommand{\Q}{\mathbb{Q}}
\providecommand{\D}{\mathbb{D}}
\providecommand{\C}{\mathbb{C}}
\providecommand{\e}{{\mathrm{e}}}
\providecommand{\vect}[1]{\overrightarrow{#1}}
\providecommand{\vectt}[1]{\overrightarrow{\mathrm{#1}}}
\providecommand{\vvt}[1]{\overrightarrow{\mathrm{#1}}}
\providecommand{\np}[2][]{{#2~#1}}
\providecommand{\vv}[1]{{\overrightarrow{#1}}}
\renewcommand{\d}{{\mathrm{d}}}
\providecommand{\degres}{{^\circ}}
\providecommand{\degre}{{^\circ}}
\providecommand{\blue}{\color{blue}}
\providecommand{\red}{\color{red}}
\providecommand{\enskip}{\quad}
\providecommand{\llbracket}{\lbrack\!\lbrack}
\providecommand{\rrbracket}{\rbrack\!\rbrack}
\providecommand{\euro}{{\text{€}}}
\providecommand{\no}{{\text{n°}}}
\providecommand{\interval}[2]{{[#1;#2]}}
\providecommand{\intervFO}[2]{{[#1;#2[}}
\providecommand{\dotfill}{\ldots}

\providecommand{\ds}{\displaystyle}
\providecommand{\textding}[1]{#1}
\providecommand{\Oij}{\(\mathrm{O};\vec{\imath},\vec{\jmath}\)}
\providecommand{\Oijk}{\(\mathrm{O};\vec{\imath},\vec{\jmath},\vec{k}\)}
\providecommand{\Ouv}{\(\mathrm{O};\vec{u},\vec{v}\)}
\providecommand{\intervFF}[2]{\ensuremath{[#1;#2]}}
\renewcommand{\euro}{\ifmmode\text{\texteuro}\else\texteuro\fi}
\DeclareUnicodeCharacter{00B0}{\ifmmode^{\circ}\else\textdegree\fi}
\DeclareUnicodeCharacter{20AC}{\ifmmode\text{\texteuro}\else\texteuro\fi}
% Énoncés : Éducation nationale. Transcription et corrigés : APMEP, Denis Vergès et contributeurs. Découpage et classement : CoopMaths, Sébastien Lozano, Éric Elter et contributeurs. Adaptation de la présentation : Maths Explorer.
% CC BY-SA : https://coopmaths.fr/www/about/
% https://coopmaths.fr/alea/static/bac/2022/tex/bac_2022_05_sujet1_asie_1_cor.tex
\begin{document}
\noindent\textit{APMEP / CoopMaths. CC BY-SA. Présentation adaptée par Maths Explorer.}\par\medskip
\medskip

\emph{Principaux domaines abordés}: Probabilités conditionnelles et indépendance. Variables aléatoires.

\medskip
















\begin{enumerate}
\item 
$R$ l'évènement \og la case obtenue est rouge\fg{} et $G$ l'évènement \og le joueur gagne la partie \fg. 
	\begin{enumerate}
		\item 
Si la case est blanche on tire 1 seul jeton ; comme il y a 3 résultats impairs sur 5 numéros on a $P_B(G) = \dfrac{3}{5} = \dfrac{6}{10} = 0,6$.
		\item 
		

Tout d'abord on a $P(B) = \dfrac{4}{12} = \dfrac13$ et donc $P(R) = \dfrac{8}{12} = \dfrac23$.

Si la case est rouge on tire successivement deux jetons : il y a $5 \times 4 = 20$ issues différentes depuis 1--2 jusqu'à 5-4 et parmi celles-ci les issues gagnantes : 1-3 ; 1-5 ; 3-1 ; 3-5 ; 5-1 et 5-3.

On a donc $P_R(G) = \dfrac{6}{20} = \dfrac{3}{10} = 0,3$ (voir l'énoncé).

On peut donc compléter l'arbre pondéré :

\begin{center}
\pstree[treemode=R,nodesepB=3pt,levelsep=3cm]{\TR{}}
{\pstree{\TR{$B$~~} \taput{$\frac13$}}
	{\TR{$G$}\taput{0,6}
	\TR{$\overline{G}$}\tbput{0,4}
	}
\pstree{\TR{$R$~~} \tbput{$\frac23$}}
	{\TR{$G$}\taput{0,3}
	\TR{$\overline{G}$}\tbput{0,7}
	}	
}
\end{center}

	\end{enumerate}

\item
	\begin{enumerate}
		\item 
D'après la loi des probabilités totales :
		
$P(G) = P(B \cap G) + P(R \cap G)$
	
$\bullet~~$$P(B \cap G) = P(B) \times P_B(G) = \dfrac13 \times 0,6 = \dfrac{0,6}{3} = 0,2$ ; 

$\bullet~~$$P(R \cap G) = P(R) \times P_R(G) = \dfrac23 \times 0,3 = \dfrac{0,6}{3} = 0,2$.

Donc $P(G) = 0,2 + 0,2 = 0,4$.
		\item 
		

Il faut trouver $P_G(B) = \dfrac{P(G \cap B)}{P(G)} = \dfrac{P(B \cap G)}{P(G)} = \dfrac{0,2}{0,4} = \dfrac2 = 0,5$.
	\end{enumerate}
\item 
$\bullet~~$$P(B) \times P(G) = \dfrac13 \times 0,4 = \dfrac{0,4}{3} = \dfrac{4}{30} = \dfrac{2}{15}$ ;

$\bullet~~$$P(B \cap G) = \dfrac13 \times 0,6 = 0,2 = \dfrac{2}{10} = \dfrac15$.

Donc $P(B) \times P(G) \ne P(B \cap G)$ : les évènements ne sont pas indépendants.
\item Un même joueur fait dix parties. Les jetons tirés sont remis dans le sac après chaque partie.

On note $X$ la variable aléatoire égale au nombre de parties gagnées.
	\begin{enumerate}
		\item 
	Les 10 épreuves sont indépendantes et à chacune la probabilité de gagner est égale à 0,4. La variable aléatoire $X$ suit donc une loi binomiale $\mathcal{B}(10~;~0,4)$.
		\item 
On a $P(X = 3 = \binom{10}{3} \times 0,4^3 \times (1 - 0,4)^7 \approx 0,215$ au millième près.
		\item 
		La calculatrice donne $P(X \leqslant 3) \approx 0,3823$, donc $P(X \geqslant 4) \approx 0,6177$, soit 0,618 au millième près.
		


Il y a plus de 6 chances sur 10 de gagner au moins 4 parties.
	\end{enumerate}
\item 
	\begin{enumerate}
		\item 
On a $p_n = P(X \geqslant 1) = 1 - P(X = 0) = 1 - 0,6^n$.
		\item 
		
Il faut trouver le plus petit entier $n$ tel que $p_n \geqslant 0,99$, soit :
		
$1 - 0,6^n \geqslant 0,99 \iff 0,01 \geqslant 0,6^n$ soit par croissance de la fonction logarithme : $\ln 0,01 \geqslant n \ln 0,6 \iff \dfrac{\ln 0,01}{\ln 0,6} \leqslant n$, car $\ln 0,6 < 0$.
		
Or $\dfrac{\ln 0,01}{\ln 0,6} \approx 9,02$.
		
Il faut donc jouer au moins 10 parties pour avoir une probabilité d'en gagner une avec une probabilité d'au moins 99\,\%.
	\end{enumerate}
\end{enumerate}

\bigskip

\end{document}
